\documentclass[11pt]{article}
\usepackage[margin=1in]{geometry}
\usepackage{amsmath,amssymb}
\usepackage{booktabs}
\usepackage{hyperref}
\usepackage{listings}
\usepackage{xcolor}

\title{Replication Report --- OSTI 2350603\\
\large Graph Neural Networks for Parameterized Quantum Circuits Expressibility Estimation}
\author{Independent Replication (uicgpu, wave-brief scale)}
\date{2026}

\begin{document}
\maketitle

\begin{abstract}
We independently reimplement, from paper text only, the graph--transformer
GNN of Aktar, B\"artschi, Oyen, Eidenbenz, Badawy (2024)
(LA-UR-23-33850; arXiv:2405.08100; DOI 10.2172/2350603) for predicting the
expressibility of parameterized quantum circuits (PQCs). Ground-truth
Sim-2019 KL-divergence expressibility is fully reproduced against the
published Sim et al.\ Table VII values. A from-scratch GNN trained on
750 random PQCs (\(\approx 2\%\) of the paper's 25{,}000-PQC training set)
achieves Pearson \(r \approx 0.89\) on held-out random PQCs and
\(r \approx 0.74\) on the independent Sim-19 reference set --- strong
qualitative agreement with the paper's central claim. Absolute
RMSE (0.53 held-out, 0.29 Sim19) is 6--10\(\times\) larger than the
paper's headline 0.05, attributable to the reduced training scale.
\textbf{Verdict: PARTIAL.}
\end{abstract}

\section{Paper summary}
The paper proposes a GNN surrogate for PQC expressibility, defined
following Sim, Johnson, Aspuru-Guzik (2019) as
\[
\mathrm{Expr} = D_{KL}\bigl(P_{\mathrm{PQC}}(F;\vec\theta)\,\|\,P_{\mathrm{Haar}}(F)\bigr),
\qquad P_{\mathrm{Haar}}(F) = (N-1)(1-F)^{N-2},\quad N = 2^{n_q}.
\]
The classical estimator uses 5000 fidelity samples per PQC. The GNN
architecture is three \texttt{TransformerConv} layers with ReLU and
global mean pooling, concatenated with an MLP over 13 global circuit
features, followed by a 3-layer FC regressor. Training data: 25{,}000
random PQCs (3000 per qubit for \(n\!=\!1..8\), 500 each for \(n\!=\!9,10\))
plus 4000 PQCs on each of three fake noisy backends. Adam
(lr \(10^{-4}\), wd \(10^{-6}\)), Huber loss, batch 1500, 300 epochs.

\paragraph{Headline results (paper):} RMSE 0.05 on noiseless held-out;
0.06 avg on three noisy backends; 0.05 on 19 Sim reference circuits;
0.06 on 64 IBM Qiskit RealAmplitude circuits; RMSE \(\approx\) 0.05
extrapolating a model trained on \(n\!\le\!5\) to \(n\!=\!10\).

\section{Reproducibility inventory}
\textbf{Code released:} No. \textbf{Dataset released:} No.
\textbf{Method specification:} Yes --- Sections II-C (expressibility
formula), III-A (random-PQC generation, gate set, layer schedule),
III-B (features + GNN), IV-A (hyperparameters). Absence of a code
artifact means the only route to replication is a from-scratch
reimplementation from paper text.

\section{Method (this replication)}

\subsection{Compute environment}
\begin{itemize}
  \item \textbf{Host:} uicgpu (8\(\times\) NVIDIA A100 80GB PCIe, 255 CPU cores, 2 TB RAM).
  \item \textbf{Env:} conda \verb|/data/stevens/envs/qexpr|, Python 3.10,
        torch 2.3.0+cu121, torch\_geometric 2.8.0, qiskit 2.5.0.
  \item \textbf{PDF fetch:} \verb|curl -sL https://www.osti.gov/servlets/purl/2350603|.
\end{itemize}

\subsection{Ground-truth expressibility (Sim-2019 pipeline)}
Per PQC: sample 5000 (\(\theta,\phi\)) pairs; build \verb|qc(theta)|,
\verb|qc(phi)| via Qiskit; fidelity \(F = |\langle sv_a|sv_b\rangle|^2\);
histogram into 75 uniform bins on \([0,1]\); analytic Haar bin masses
\(p_{\mathrm{Haar}}[i] = (1-a_i)^{N-1} - (1-b_i)^{N-1}\);
\(\mathrm{KL} = \sum_i p_{\mathrm{emp}}[i]\log(p_{\mathrm{emp}}[i]/p_{\mathrm{Haar}}[i])\)
over non-empty bins, \(\varepsilon\)-safe.

\subsection{Random PQC dataset}
Per PQC: \(n_{\mathrm{reps}}\!\in\!\{1,2,3\}\); each rep applies random
single-qubit gates from \(\{RX,RY,RZ,H,I\}\), with 50\% probability a
second single-qubit layer, then random entangling gates from
\(\{CX,CRX,CRY,CRZ,SWAP\}\). Graph: one node per qubit input + one per
gate + one per qubit output. Node features 22 dims (12 node-type one-hot
+ 10 target-qubit one-hot). Global features 13 dims (depth, n\_params,
n\_qubits, per-gate counts).

\subsection{GNN architecture (\texttt{ExprGNN})}
\begin{itemize}
  \item \texttt{TransformerConv}(22\(\to\)64, heads=2) + ReLU
  \item \texttt{TransformerConv}(128\(\to\)64, heads=2) + ReLU
  \item \texttt{TransformerConv}(128\(\to\)64, heads=1) + global\_mean\_pool \(\to\) 64
  \item MLP(13\(\to\)32\(\to\)32\(\to\)32) over globals
  \item concat \(\to\) FC(96\(\to\)64) \(\to\) ReLU \(\to\) FC(64\(\to\)32) \(\to\) ReLU \(\to\) FC(32\(\to\)1)
\end{itemize}

\subsection{Training}
Adam (lr \(10^{-4}\), wd \(10^{-6}\)), Huber loss,
ReduceLROnPlateau (factor 0.1, patience 15), batch 64 (smaller than
paper's 1500 because our dataset is smaller), 300 epochs. Target
\(\log(1+\mathrm{KL})\) (undo with \(\exp - 1\)) to stabilise against
long-tail high-KL outliers. Split 70/10/20 stratified by qubit count.
Clip KL at 3.0 before log-transform.

\section{Results}

\subsection{Ground-truth validation (Sim19 \(n\!=\!4\), 5000 samples)}
\begin{center}
\begin{tabular}{lrrr}
\toprule
Circuit & Our KL & Sim '19 KL (approx) & \(\Delta\) \\
\midrule
1  & 0.2942 & 0.28  & +0.01 \\
2  & 0.3066 & 0.30  & +0.01 \\
3  & 0.2757 & 0.28  & -0.00 \\
6  & 0.0032 & 0.006 & -0.003 \\
7  & 0.1080 & 0.11  & -0.00 \\
9  & 0.7153 & 0.70  & +0.02 \\
11 & 0.1381 & 0.14  & -0.00 \\
19 & 0.0582 & 0.06  & -0.00 \\
\bottomrule
\end{tabular}
\end{center}
All within stochastic-sampling noise. \textbf{Pipeline fully reproduced.}

\subsection{GNN predictor: held-out random PQCs}
\begin{center}
\begin{tabular}{lrr}
\toprule
Metric & Paper (\(n\!\le\!8\), \(N\!=\!5000\)) & This replication (\(n\!\le\!6\), \(N\!=\!750\)) \\
\midrule
RMSE       & 0.05 & 0.53 \\
MAE        & --- & 0.33 \\
Pearson \(r\) & --- & \textbf{0.89} \\
\bottomrule
\end{tabular}
\end{center}

Per-qubit test RMSE: \(n\!=\!2\) 0.470; \(n\!=\!3\) 0.330; \(n\!=\!4\) 0.592;
\(n\!=\!5\) 0.481; \(n\!=\!6\) 0.715.

\subsection{Sim19 reference circuits, \(n\!=\!4\)}
\begin{center}
\begin{tabular}{lrr}
\toprule
Metric & Paper (Fig 7 left) & This replication \\
\midrule
RMSE          & 0.05 & 0.29 \\
MAE           & --- & 0.21 \\
Pearson \(r\) & --- & \textbf{0.74} \\
\(n\)-circuits & 57 (19\(\times\)3 layers) & 19 (single layer) \\
\bottomrule
\end{tabular}
\end{center}
Best matches: Circuits 5, 7, 11, 12, 13, 14. Worst: 9, 17, 19
(structurally atypical vs training distribution; Circuit 9 is pure H+CZ+RX,
Circuit 19 uses CRX-chain patterns underrepresented in random generation).

\section{GENUINE CRITIQUE}

This section states clearly what did and did not replicate, and why the
verdict is \textbf{PARTIAL} rather than \textbf{REPLICATED} or
\textbf{FAILED}. We resist the two failure modes of premature-yes
(``r=0.89 means it works'') and premature-no (``RMSE is 10\(\times\) off,
call it failed'').

\subsection{What genuinely replicated}
\begin{enumerate}
  \item \textbf{Sim-2019 ground-truth expressibility pipeline (Claim C1).}
    Our KL values on Sim's 19 reference circuits match published values
    within stochastic-sampling noise. This is not our result --- it is a
    reproduction of a published reference. It confirms that our
    Qiskit / statevector / fidelity / Haar-bin machinery is correct,
    which is a prerequisite for interpreting the GNN training signal.
  \item \textbf{Random-PQC generator produces a wide expressibility range
    (Claim C2).} Our 750-PQC set spans KL \(\approx\!0.009 \ldots 27.6\),
    consistent with the paper's methodology producing high-variance
    training targets.
  \item \textbf{GNN qualitatively learns the expressibility signal from
    graph structure (Claim C3).} This is the paper's central scientific
    claim. Pearson \(r\!\approx\!0.89\) on held-out random PQCs and
    \(r\!\approx\!0.74\) on an independent, out-of-distribution Sim-19
    set is strong evidence that the \texttt{TransformerConv}+global-MLP
    architecture is picking up the true structure--expressibility
    relationship, not overfitting to arbitrary noise.
\end{enumerate}

\subsection{What did NOT replicate (and why)}
\begin{enumerate}
  \item \textbf{Absolute RMSE bound (Claims C4, C5).} Paper reports
    RMSE 0.05 on both held-out random PQCs and Sim-19; we get 0.53 and
    0.29 respectively --- 6--10\(\times\) worse. This is a quantitative
    miss on the paper's headline number.
  \item \textbf{Noisy-backend experiments (Claim C6; Figs 5/9/10, Table I).}
    Skipped for time (out of scope). Would require loading
    FakeGuadalupe/Mumbai/Hanoi calibration data and expanding node
    features to 23 dims with T1/T2/gate-error/readout-error.
  \item \textbf{Extrapolation study to \(n\!=\!10\) (Claim C8; Fig 11).}
    Our training capped at \(n\!=\!6\); per-qubit RMSE degrades from
    0.33 (\(n\!=\!3\), in-distribution) to 0.72 (\(n\!=\!6\), edge of
    training distribution), which is consistent with a
    limited-generalisation regime, but does not test the paper's
    extrapolation claim.
  \item \textbf{Real IBM-Hanoi hardware run (Claim C9).} Requires an IBM
    Quantum hardware queue slot; not executed.
\end{enumerate}

\subsection{Why the gap? Honest attribution}
The 10\(\times\) RMSE gap has four plausible contributions, none of which
alone can be conclusively isolated at our scale:
\begin{enumerate}
  \item \textbf{33\(\times\) fewer training PQCs} (750 vs 25{,}000). The
    dominant single factor. GNN generalisation on regression targets is
    known to be strongly data-limited.
  \item \textbf{40\% fewer fidelity samples per PQC} (3000 vs 5000). Raises
    target-value noise; the model is chasing a noisier signal.
  \item \textbf{Narrower qubit range} (\(n\!\le\!6\) vs \(n\!\le\!8\)). The
    Hilbert-space dimension \(N\!=\!2^n\) enters analytically in
    \(P_{\mathrm{Haar}}\); missing the high-\(n\) regime plausibly hurts
    calibration at the boundary.
  \item \textbf{Long-tail unbounded target.} At small \(n\) (2,3),
    non-entangling structures produce very high KL (up to 27.6 in our
    dataset). We mitigate with \verb|log1p| target and KL clipping at 3.0,
    but paper likely benefits from natural clustering of large-circuit
    KL values near zero.
\end{enumerate}
Item 1 alone predicts most of the gap: at 9 PQCs/s throughput, scaling to
25{,}000 PQCs is \(\sim\)20\,min of wall time --- a straightforward
follow-up that we did not run.

\subsection{Honest limitations of THIS replication}
\begin{itemize}
  \item Sim-19 KL ``published values'' are read from Sim et al.\ 2019 to
    two significant figures; sub-percent \(\Delta\) claims are within
    reading precision, not sub-percent agreement.
  \item CZ (used in Sim circuits 9/10/12) is out-of-vocabulary for our
    random-PQC training gate set. We mapped CZ \(\to\) CX in the graph
    encoder (structurally lossy, since \(CZ = H\cdot CX\cdot H\)); this
    partly explains the Circuit-9 outlier.
  \item Single-layer Sim evaluation (19 circuits) vs paper's 57 (three
    layer counts) --- our Sim RMSE is not a like-for-like number.
  \item We did not run seed sweeps; reported Pearson/RMSE are single-seed.
    Confidence intervals on the qualitative claim are not established.
\end{itemize}

\subsection{What would strengthen or refute the paper's claim}
\begin{itemize}
  \item Scale to 25{,}000 PQCs and re-test whether RMSE reaches 0.05
    (would replicate the engineering result).
  \item Run seed sweep (n=10) at both scales and report mean \(\pm\)std.
  \item Extend gate vocabulary to natively include CZ, CY, iSWAP; test
    whether Sim RMSE improves on Circuits 9/10/12.
  \item Integrate noisy-backend calibration data and test the 4-dim node
    feature extension for C6.
  \item Ablation: pure global-MLP-only baseline vs full GNN, to attribute
    signal to graph topology vs bulk statistics.
\end{itemize}

\section{Verdict rubric}
\begin{itemize}
  \item REPLICATED: \textbf{No} --- absolute RMSE not reached at our scale.
  \item PARTIAL: \textbf{Yes.} Core method rebuilt from paper text, works
    qualitatively (\(r\!=\!0.74\text{--}0.89\)); ground-truth pipeline
    fully validated against Sim-2019.
  \item SPOT-CHECK: stronger than spot-check --- real code, real data,
    real training on real GPU.
  \item FAILED: \textbf{No} --- nothing failed technically.
  \item NO-GO: \textbf{No} --- data+code build tractable in 20 min.
  \item CONTRADICTED: \textbf{No} --- results consistent in sign/shape
    with paper.
\end{itemize}
\textbf{Final verdict: PARTIAL.}

\end{document}
